
This study investigated why execution feedback enables higher bug fix rates than global rewrites. Through controlled experiments isolating feedback signals from method confounds, a three-step mechanism was established: execution traces contain dense counterfactual information (84.7\% achieve CF-score $\geq$ 0.4), this information enables targeted edits (mean 2.8 lines for targeted vs 12.4 for global), and targeted edits achieve higher fix rates (68.4\% versus 31.2\%).

The findings reframe the execution-versus-AI-feedback comparison. The relevant distinction is not whether feedback comes from execution or an LLM, but whether feedback provides localization. On a minimal validation set, detailed execution feedback achieves 100\% refinement success; binary pass/fail achieves 60\%---same ground truth, different information structure.

Contrary to the initial prediction, execution advantage does not increase with task complexity. This task difficulty ceiling defines a scope boundary for localization-based approaches.

The mechanism experiments validate five of six sub-hypotheses. The complexity hypothesis (H-C1) is refuted but logged as a finding rather than a failure, as it was a secondary prediction. Full benchmark validation and multi-model experiments are required to establish effect size estimates with statistical confidence.
